\documentclass[12pt, a4paper]{article}
\usepackage{graphicx} % Required for inserting images
\graphicspath{{images/}{images-desafios/}}
\usepackage[a4paper, total={7.5in, 10in}]{geometry}
\usepackage{float} % pra conseguir usar o \begin{figure}[H]
\usepackage{indentfirst} % pra fazer tab n 1 paragrafos dos capitulos
% package pra tabela
\usepackage{multirow}
\usepackage{array}

%pacote pra codigo
\usepackage{listings} 
\usepackage{xcolor}
\definecolor{codegreen}{rgb}{0,0.6,0}
\definecolor{codegray}{rgb}{0.5,0.5,0.5}
\definecolor{codepurple}{rgb}{0.58,0,0.82}
\definecolor{backcolour}{rgb}{0.95,0.95,0.92}

\lstdefinestyle{myCodeStyle}{
	backgroundcolor=\color{backcolour},   
	commentstyle=\color{codegreen},
	keywordstyle=\color{cyan},
	numberstyle=\tiny\color{codegray},
	stringstyle=\color{codepurple},
	basicstyle=\ttfamily\footnotesize,
	breakatwhitespace=false,         
	breaklines=true,                 
	captionpos=b,                    
	keepspaces=true,                 
	%numbers=left,                    
	numbersep=5pt,                  
	showspaces=false,                
	showstringspaces=false,
	showtabs=false,                  
	tabsize=2
}
\lstset{style=myCodeStyle}
\renewcommand{\lstlistingname}{Código}%Muda o termo "Listing":  Listing -> Código
\renewcommand{\lstlistlistingname}{Lista de \lstlistingname s}% List of Listings -> Lista de Códigos


\newcommand{\unit}[1]{\ensuremath{\, \mathrm{#1}}} % comando pra tirar do math mode a unidade de medida

\def\tamanhoGrafico{0.8}

\newcommand{\ohm}{\unit{\Omega} }
\newcommand{\kohm}{\unit{k\Omega}}

\title{
	Laboratório Digital A - PCS3335
	\\ Proposta de Projeto
}

\author{
	Alunos:
	\\ Rafael Hipólito Rodrigues - N.USP: 14604308
	\\ {{PERSON-2}} - N.USP: {{PERSON-ID-2}}
	\\ \\
	Turma 3 - Bancada A6
	\\ Professor: Bruno de Carvalho Albertini
}

%\date{}

\begin{document}
	
	\maketitle
	
	\section*{Proposta de Projeto: Calculadora com Notação Polonesa Reversa e memória} 
	
	A proposta do grupo é elaborar uma calculador simples que utilizar Notação Polonesa Reversa para realizar operações 
	sobre valores que estão armazenados na memória. A calculadora realizaria as 4 básicas operações aritméticas: soma, subtração, 
	multiplicação e divisão. A ideia é utilizar a memoria como uma pilha de números. Cada numero novo acrescentado na memoria seria 
	adicionado no topo da pilha. As operações seriam realizadas entre os 2 números no topo da pilha, esses 2 números seriam retirados
	da pilha e o resultado seria armazenado no topo da pilha.
	
	
	
	\section*{Modos de funcionamento}
	
	A calculadora terá 3 modos de funcionamento:
	
	\subsection*{Modo 1: armazenamento de números na memoria como uma pilha}
	No 1° modo o usuário apenas armazenaria números na memoria, como se estivesse adicionando valores ao topo de uma pilha.
	
	\subsection*{Modo 2: realização de operações aritméticas com os valores da pilha}
	No 2° modo o usuário apenas realizaria operações aritméticas com os valores que já estiverem na armazenados na memoria. Ao selecionar 
	uma operação aritméticas, a calculadora pegaria os 2 números mais ao topo da pilha, realizaria a operação entre eles, removeria esses 2 números da 
	pilha, e armazenaria no topo da pilha o resultado da operação.
	
	\subsection*{Modo 3: exibição dos valores da pilha}
	No 3° modo o usuário poderá ler os valores que ja estão armazenados na memoria.
	
	
	
	\section*{Circuitos combinatórios}
	A entidade da calculadora devera ter circuitos combinatórios para cada operações aritméticas realizada e outros necessários para implementar a 
	logica de funcionamento de cada modo. 
	
	\section*{Máquina de estados}
	A entidade devera ter uma máquina de estados para cada modo de funcionamento. 
	
	\section*{Elemento de memória}
	A entidade terá um elemento de memoria de conteúdo modificável e de acesso aleatório (similar a memoria RAM implementada em SD2).
	
	
	
	\section*{Características da calculadora}
	Algumas características e detalhes da calculadora:
	
	\begin{itemize}
		\item O valores serão utilizados em hexadecimal. O grupo escolheu hexadecimal pois será mais fácil de converter para binário e será possível
		exibir valores maiores nos 6 displays de 7 segmentos
		
		\item O valores será armazenados em 21 bits: 1 bits para o sinal e 20 bits para o valor numérico. Assim o valor máximo será FFFFF em 
		hex (1111 1111 1111 1111 1111 em bin). A escolha dessa quantidade é para conseguir exibir o valores no 6 displays de 7 segmentos. Sendo os 5
		primeiros displays para o valor numérico e o 6° display para o sinal. O sinal no valor binário será utilizara logica de complemento de 2.
		
		\item Caso haja overflow nas operações, um LED será acesso para indicar erro de overflow, contudo o valor que sofreu overflow será armazenado 
		normalmente na memoria.
		
		\item Os valores utilizados serão apenas números inteiros com sinal.
		
		\item A operação de divisão resultara no valor inteiro da divisão.
		
		\item Caso haja uma divisão por zero, um LED será acesso para indicar erro de divisão por zero, e a operação não será realizada.
		
		\item O input dos dados no Modo 1 será utilizando um teclado de membrana matricial 4x4 disponibilizado pelo Lab. Dig.. Também será utilizados 
		as chaves da FPGA ou as GPIOs para fazer o controle entre os modos de funcionamento.
		
		\item No modo 1 apenas será possível digitar um numero de 5 dígitos hexadecimal e um sinal para o numero. 
		
		\item Caso haja apenas 1 numero armazenado na memoria, não será permitido realizar uma operação aritmética no modo 2.
		
	\end{itemize}
	
	OBS.: A dupla decidiu essas características para o projeto, porém pode correr da dupla decidir modificar algumas delas a fim de melhorar o 
	resultado final do projeto.
	
	
	
	\section*{Cronograma de execução para 3 experiências (9, 10 e 11)}
	
	\begin{itemize}
		\item Experiência 9 - 20/05: elaboração dos circuitos combinatórios de operações aritméticas; elaboração da entidade de memoria; 
		elaboração de logica que utiliza o teclado de membrana matricial para o input de números.
		
		\item Experiência 10 - 27/05: implementação dos 3 modos de funcionamento.
		
		\item Experiência 11 - 03/06: finalização e apresentação da calculadora.
	\end{itemize}
	
	
\end{document}
